\section{MLIR 与渐进式降低}
\label{sec:mlir}

LLVM IR 适合表达显式控制流、内存和机器式标量操作，但领域操作若在优化前被拆解，其高层结构可能难以恢复。MLIR 通过可扩展方言（dialect）和统一的操作、区域与类型基础设施容纳多个抽象层级\cite{lattner-et-al-mlir-2021}。这一设计与前述案例的表示演化具有共同问题：某类信息应保留至哪些变换完成，以及转换边界如何验证。

对阶乘循环，一条概念性路径可从 \code{func}、\code{arith} 与 \code{scf} 方言开始，其中函数、整数运算和结构化循环仍然显式；降低至 \code{cf} 后，循环成为基本块、分支和块参数；随后转换至 LLVM 方言并翻译为 LLVM IR，继续使用 LLVM 优化与 RV64 后端。方言转换框架通过转换目标声明操作合法性，以重写模式替换非法操作，并使用类型转换器处理混合表示间的类型衔接\cite{mlir-dialect-conversion-2026,mlir-llvm-target-2026}。Toy 教程中的部分降低展示了多种方言在转换期间共存的方式\cite{mlir-toy-lowering-2026}。

异构计算进一步强化了多级表示的需求。AscendNPU IR 官方 VecAdd 材料展示了从上层张量/计算描述向设备相关表示逐步降低的设计背景\cite{ascend-vecadd-2026,ascend-architecture-2026}。该外部案例说明，领域结构、存储层次和硬件映射可能需要在不同阶段分别显式化；它不构成本文环境中的 Ascend 编译或设备执行证据。

渐进式降低的工程代价包括方言合法性规则、重写模式、类型桥接、诊断以及版本维护。多级表示适用于确有领域或硬件结构需要延迟降低的系统；对于本研究中的小型标量程序，LLVM IR 已足以承载所需分析和代码生成。由此可见，表示层级的数量应由待保留信息及其变换需求决定。
